\PassOptionsToPackage{expansion=false}{microtype}
\documentclass[bachelor, och, coursework]{SCWorks}
\usepackage[utf8]{inputenc}
\usepackage[T2A]{fontenc}
\usepackage{graphicx}
\usepackage[sort,compress]{cite}
\usepackage{amsmath}
\usepackage{amssymb}
\usepackage{amsthm}
\usepackage{fancyvrb}
\usepackage{longtable}
\usepackage{array}
\usepackage{multirow}
\usepackage[english,russian]{babel}
\usepackage{caption}
\captionsetup[figure]{font=normalsize, labelfont=normalsize}
\usepackage[hidelinks]{hyperref}
\sloppy
\usepackage{listings}
\usepackage{xcolor}
\usepackage{float}
\usepackage{microtype}
\usepackage{url}

\lstset{
	basicstyle=\small\ttfamily,
	numbers=left,
	stepnumber=1,
	numbersep=8pt,
	frame=none,
	xleftmargin=2em,
	xrightmargin=0.5em,
	breaklines=true,
	breakatwhitespace=false,
	showstringspaces=false,
	tabsize=4,
	captionpos=b,
	language=Python,
	inputencoding=utf8,
	extendedchars=true,
	aboveskip=\medskipamount,
	belowskip=\medskipamount,
}

\begin{document}
	
	\chair{дискретной математики и информационных технологий}
	\title{Сравнение методов квантования LLM (INT8, INT4, NF4)}
	\course{2}
	\group{221}
	\napravlenie{09.03.01 -- Информатика и вычислительная техника}
	\author{Соколов Даниил Васильевич}
	\chtitle{к.\,ф.-м.\,н., доцент}
	\chname{Л.\,Б.\,Тяпаев}
	\satitle{ст. преподаватель}
	\saname{П.\,О.\,Дмитриев}
	\date{2026}
	
	\maketitle
	\tableofcontents
	
	\abbreviations
	\begin{description}
		\item[LLM] Large Language Models, большие языковые модели
		\item[INT8] 8-битное целочисленное квантование
		\item[INT4] 4-битное целочисленное квантование
		\item[NF4] 4-bit NormalFloat, квантование с нормальным распределением уровней
		\item[GPTQ] Generative Pre-trained Transformer Quantization
		\item[QLoRA] Quantized Low-Rank Adaption
		\item[MMLU] Massive Multitask Language Understanding
		\item[MSQE] Mean Squared Quantization Error, среднеквадратическая ошибка квантования
		\item[FP16] 16-битное число с плавающей точкой
		\item[SQNR] Signal-to-Quantization-Noise Ratio
	\end{description}
	
	\intro
	
	Число параметров больших языковых моделей (Large Language Models, LLM) быстро увеличивается, что повышает требования к вычислительным ресурсам при обучении и инференсе. Модели, содержащие миллиарды параметров, показывают высокие результаты в решении разнообразных задач обработки естественного языка, однако их практическое применение ограничивается высокими требованиями к оперативной памяти и вычислительным ресурсам. В этом контексте методы квантования являются важным направлением оптимизации, позволяющим снизить разрядность представления весовых коэффициентов и активаций при сохранении приемлемого уровня точности модели.
	
	Актуальность исследования методов квантования обусловлена несколькими факторами. Во-первых, необходимость развёртывания LLM на устройствах с ограниченными ресурсами требует разработки эффективных методов сжатия моделей. Во-вторых, экономическая эффективность применения LLM в промышленности связана с оптимизацией вычислительных затрат. В-третьих, энергопотребление крупных нейросетевых систем является значимым фактором при их промышленном внедрении.
	
	Методы квантования нейронных сетей исследуются с начала 2010-х годов. Первые работы были посвящены в основном свёрточным архитектурам для задач компьютерного зрения. Однако специфика трансформерных архитектур, лежащих в основе современных LLM, требует адаптации классических подходов к сжатию моделей. Распределение весовых коэффициентов в языковых моделях характеризуется наличием выбросов и несимметричностью, что существенно усложняет применение традиционных методов равномерного квантования.
	
	Несмотря на активные исследования последних лет, некоторые вопросы квантования языковых моделей остаются недостаточно изученными. К дискуссионным вопросам относятся оптимальный баланс между степенью сжатия модели и деградацией ее производительности, выбор оптимальной гранулярности квантования, а также влияние различных методов на специфические аспекты поведения модели, такие как способность к рассуждениям и генерации длинных контекстов. Требуется разработка унифицированной методологии выбора метода квантования в зависимости от целевой задачи и аппаратных ограничений.
	
	\textbf{Цель курсовой работы} состоит в проведении сравнительного анализа методов квантования больших языковых моделей INT8, INT4 и NF4 с точки зрения их влияния на качество генерации, вычислительную эффективность и применимость к различным типам задач.
	
	Для достижения поставленной цели необходимо решить следующие \textbf{задачи}:
	\begin{itemize}
		\item исследовать математические принципы сокращения разрядности параметров в глубоких нейронных сетях и формализовать типы ошибок дискретизации;
		\item провести сравнительный анализ алгоритмов целочисленного квантования от симметричного до поблочного подходов;
		\item систематизировать существующие экспериментальные результаты применения методов INT8, INT4 и NF4 к современным языковым моделям;
		\item проанализировать влияние низкоразрядного квантования на ключевые метрики качества и сформулировать рекомендации по выбору метода.
	\end{itemize}
	
	\textbf{Объектом исследования} выступают большие языковые модели трансформерной архитектуры и методы их оптимизации.
	
	\textbf{Предметом исследования} являются методы квантования весовых коэффициентов LLM с различной разрядностью представления: INT8, INT4 и NF4.
	
	\textbf{Теоретическую базу исследования} составляют фундаментальные концепции теории информации, численных методов и оптимизации нейронных сетей. Математический аппарат работы опирается на теорию квантования сигналов, статистическую теорию оценивания и линейную алгебру. Методологической основой служат работы по архитектуре трансформеров, механизмам внимания и методам компрессии нейросетевых моделей.
	
	\textbf{Методы исследования} включают аналитические методы исследования алгоритмов квантования, численное моделирование процессов дискретизации весовых коэффициентов, сравнительный анализ экспериментальных результатов из научных публикаций и технических отчетов, статистические методы обработки данных бенчмарков. Материалы исследования охватывают результаты применения различных методов квантования к семействам моделей LLaMA, GPT, BLOOM и OPT различного размера. Использованы данные открытых бенчмарков GLUE, SuperGLUE, MMLU, а также специализированные наборы данных для оценки перплексии на корпусах WikiText и C4.
	
	\textbf{Научная новизна работы} состоит в систематизации современных методов квантования LLM, включая относительно новый метод NF4, основанный на неравномерном распределении квантовальных уровней, оптимизированном для нормального распределения весов.
	
	\textbf{Проблема}, решаемая в работе, связана с ростом размера LLM. Например, модель LLaMA-7B в формате FP16 занимает около 13 ГБ видеопамяти. Это превышает возможности многих потребительских GPU (NVIDIA RTX 3060 — 12 ГБ, RTX 4060 Ti — 8 ГБ). В результате стандартный FP16-инференс на таком оборудовании невозможен.
	
	Квантование позволяет уменьшить размер модели, но ценой возможной деградации качества. Возникает вопрос: какой метод квантования (INT8, INT4 или NF4) обеспечивает оптимальный баланс между размером модели, скоростью инференса и точностью для различных сценариев использования?
	
	В работе проводится сравнение этих методов по следующим критериям: итоговый размер модели (ГБ); скорость инференса (токен/с); деградация качества (перплексия, точность на бенчмарке MMLU). Результаты позволяют дать практические рекомендации: для GPU с 12 ГБ памяти подходит INT8-квантование, для 6–8 ГБ — INT4 GPTQ, для задач дообучения на ограниченном оборудовании — NF4 + QLoRA.
	
	\section{Теоретические принципы сокращения разрядности параметров в глубоких нейронных сетях}
	
	\subsection{Математическая формализация процесса квантования и анализ типов ошибок дискретизации}
	
	Квантование — это отображение множества вещественных чисел (например, FP32) на конечное дискретное множество квантовальных уровней. В контексте нейронных сетей квантование применяется к весовым коэффициентам, активациям или градиентам с целью снижения требований к памяти и вычислительным ресурсам. Математически процесс квантования формализуется как функция \(Q: \mathbb{R} \to \mathcal{Q}\), где \(\mathbb{R}\) обозначает множество вещественных чисел (или его подмножество с конечной точностью, например, FP32 или FP16), а \(\mathcal{Q}\) представляет собой конечное дискретное множество квантовальных уровней.
	
	Для унифицированного представления различных методов квантования используется следующая общая формула:
	\[
	Q(x) = s \cdot \operatorname{round}\left(\frac{x - z}{s}\right) + z \tag{1}
	\]
	где \(x\) — исходное вещественное значение; \(s\) — масштабный коэффициент (scale factor); \(z\) — параметр смещения (zero-point); \(\operatorname{round}(\cdot)\) — операция округления до ближайшего целого числа. Параметры \(s\) и \(z\) определяют характеристики квантования и вычисляются на основе статистических свойств квантуемых данных. Таким образом, из (1) следует, что точность квантования определяется выбором масштабного коэффициента \(s\) и параметра смещения \(z\).
	
	Для симметричного квантования, когда диапазон значений центрирован относительно нуля, параметр смещения \(z\) принимается равным нулю, что даёт формулу \(Q(x) = s \cdot \operatorname{round}(x / s)\). В этом случае масштабный коэффициент определяется как \(s = (x_{\max} - x_{\min}) / (2^{b} - 1)\), где \(b\) обозначает разрядность квантования в битах, \(x_{\max}\) и \(x_{\min}\) — максимальное и минимальное значения в квантуемом диапазоне \cite{llmint8}.
	
	Если распределение значений несимметрично относительно нуля, используется асимметричное квантование с параметром смещения \(z \neq 0\) (формула 1). Таким образом, выражение (1) является универсальным и описывает оба случая. Во всех практических сценариях масштабный коэффициент \(s > 0\), а результат округления ограничен диапазоном целевого целочисленного типа (например, int8 или uint8).
	
	Асимметричное квантование применяется в случаях, когда распределение значений существенно смещено относительно нуля. Параметр смещения \(z\) вычисляется как \(z = -\operatorname{round}(x_{\min} / s)\), что обеспечивает полное использование доступного диапазона квантовальных уровней. Асимметричный подход особенно эффективен для квантования активаций, которые часто характеризуются несимметричным распределением, например, после применения функции ReLU \cite{awq}.
	
	Ошибка квантования \(\varepsilon(x)\) определяется как разность между исходным значением и его квантованным представлением: \(\varepsilon(x) = x - Q(x)\). Величина этой ошибки напрямую влияет на точность вычислений в квантованной нейронной сети. Для равномерного квантования с шагом квантования \(\Delta = s\) максимальная абсолютная ошибка ограничена величиной \(\Delta/2\), что следует из свойств операции округления. Следовательно, ошибка квантования ограничена половиной шага квантования \(\Delta/2\), что определяет максимальную потерю точности при преобразовании.
	
	Среднеквадратическая ошибка квантования (Mean Squared Quantization Error, MSQE) служит важной метрикой для оценки качества квантования и определяется выражением:
	\[
	\operatorname{MSQE} = \mathbb{E}\left[(x - Q(x))^{2}\right] \tag{2}
	\]
	где \(\mathbb{E}[\cdot]\) обозначает математическое ожидание. Для равномерного распределения исходных значений в пределах диапазона квантования среднеквадратическая ошибка составляет \(\operatorname{MSQE} = \Delta^{2}/12\), что демонстрирует квадратичную зависимость ошибки от шага квантования.
	
	Отношение сигнал/шум квантования (Signal-to-Quantization-Noise Ratio, SQNR) характеризует относительное качество квантованного представления:
	\[
	\operatorname{SQNR} = 10 \cdot \log_{10}\left(\frac{\sigma_{x}^{2}}{\operatorname{MSQE}}\right) \tag{3}
	\]
	где \(\sigma_{x}^{2}\) — дисперсия исходного сигнала. Для \(b\)-битного равномерного квантования теоретическое значение SQNR приближенно составляет \(6.02b + 1.76\) дБ, что указывает на приблизительно 6 дБ улучшения на каждый дополнительный бит разрядности \cite{nf4}. Из (3) видно, что каждый дополнительный бит разрядности улучшает соотношение сигнал/шум примерно на 6 дБ, что количественно обосновывает выигрыш от увеличения разрядности.
	
	Калибровка квантования представляет собой важный этап, определяющий параметры \(s\) и \(z\) на основе статистических характеристик данных. Существуют различные подходы к калибровке.
	
	\textbf{Метод min-max} определяет диапазон квантования на основе минимального и максимального наблюдаемых значений. Однако этот подход чувствителен к выбросам в данных, что может приводить к неэффективному использованию квантовальных уровней для основной массы значений.
	
	\textbf{Процентильный метод} использует заданные перцентили распределения для определения диапазона квантования, игнорируя экстремальные выбросы. Типично используются перцентили 0.1\% и 99.9\%, что обеспечивает робастность к редким аномальным значениям.
	
	\textbf{Метод минимизации Kullback-Leibler дивергенции} выбирает параметры квантования, минимизирующие расхождение между распределением исходных и квантованных активаций, что теоретически оптимизирует сохранение информации \cite{smoothquant}.
	
	Выбор метода калибровки зависит от распределения данных: для данных с выбросами предпочтителен процентильный метод, для нормальных распределений — метод min-max.
	
	\subsection{Гранулярность квантования и её влияние на точность}
	
	Гранулярность квантования определяет уровень, на котором применяются единые параметры квантования.
	
	\textbf{Послойное квантование} (per-layer) использует одни параметры \(s\) и \(z\) для всего слоя нейронной сети. Это самый простой подход с минимальными накладными расходами, но он плохо адаптируется к случаям, когда разные каналы слоя имеют сильно различающиеся масштабы значений.
	
	\textbf{Поканальное квантование} (per-channel) применяет отдельные параметры для каждого выходного канала. Для весовой матрицы \(W\) размерности \([n, m]\) определяются \(n\) пар параметров \((s_i, z_i)\), где \(i = 1, \dots, n\). Этот подход учитывает существенную вариативность масштабов активаций различных нейронов, характерную для глубоких сетей \cite{microscaling}. Эмпирические исследования показывают, что поканальное квантование позволяет достичь существенно меньшей деградации точности по сравнению с послойным подходом, особенно при агрессивном снижении разрядности до 4 бит и ниже \cite{gptq}.
	
	\textbf{Групповое квантование} (group-wise) занимает промежуточное положение. Весовые коэффициенты делятся на группы заданного размера \(G\) (например, 32, 64 или 128), и для каждой группы определяются индивидуальные параметры квантования. Для одномерного весового вектора \(w\) длины \(N\) групповое квантование разбивает его на \(\lceil N / G \rceil\) групп. Для \(k\)-той группы вычисляются параметры \(s_k\) и \(z_k\) на основе статистики значений, попадающих в эту группу.
	
	Групповое квантование адаптируется к локальным вариациям параметров, требуя при этом умеренных затрат на хранение дополнительных параметров. Для весовой матрицы размера \([n, m]\) при размере группы \(G\) требуется хранить приблизительно \((n \cdot m) / G\) пар параметров \((s, z)\), что при типичном \(G = 64\) составляет около 3\% дополнительных данных при использовании 16-битного представления параметров.
	
	\subsection{Продвинутые методы квантования}
	
	\textbf{Адаптивное округление (AdaRound)} — метод оптимизации квантования весов путём тонкой настройки операции округления. Вместо детерминистического округления к ближайшему целому метод вводит параметр \(v \in [0, 1]\) для каждого веса, определяющий направление округления: \(Q(w) = s \cdot (\lfloor w / s \rfloor + h(v))\), где \(\lfloor \cdot \rfloor\) — операция округления вниз, а \(h(v)\) — сигмоидальная функция, приближающая индикаторную функцию округления вверх \cite{sparsegpt}. AdaRound позволяет уменьшить ошибку квантования за счёт оптимизации направления округления, но требует дополнительных вычислительных затрат на калибровку.
	
	Параметры \(v\) оптимизируются путем минимизации расхождения между выходами слоя в исходной и квантованной версиях на калибровочном наборе данных. Целевая функция формулируется как \(L = \|Wx - Q(W)x\|^{2}\), где \(x\) представляет активации с предыдущего слоя. Оптимизация выполняется методом градиентного спуска, причем градиенты вычисляются относительно параметров \(v\) с использованием прямого пропуска через \(h(v)\). После оптимизации значения \(v\) бинаризуются: округление вверх применяется при \(v > 0.5\), иначе округление вниз.
	
	\textbf{Метод GPTQ (Generative Pre-trained Transformer Quantization)} специализирован для квантования больших языковых моделей и основан на оптимизации квантования с учётом структуры весовой матрицы и статистики активаций. GPTQ последовательно квантует столбцы весовой матрицы, компенсируя ошибку квантования каждого столбца путём корректировки ещё неквантованных столбцов. Компенсация основана на обратной матрице Гессе функции потерь, вычисленной по калибровочным данным \cite{gptq}.
	
	Алгоритм GPTQ для весовой матрицы \(W\) размерности \([n, m]\) включает следующие этапы. Вычисляется матрица Гессе \(H = 2 X X^{T} / N\), где \(X\) — матрица активаций калибровочного набора размера \([n, N]\). Для каждого столбца \(j = 1, \dots, m\): квантуется столбец \(W_{:,j}\), вычисляется ошибка \(e_j = W_{:,j} - Q(W_{:,j})\), корректируются оставшиеся столбцы \(W_{:,k}\) для \(k > j\) путем вычитания \(H^{-1}_{:,j} e_j H_{j,k} / H_{j,j}\). Метод обеспечивает высокое качество квантования при разрядности 4 бита и ниже для моделей масштаба десятков миллиардов параметров.
	
	\subsection{Аппаратная поддержка квантования}
	
	Аппаратная поддержка различных типов квантования варьируется в зависимости от платформы. Современные GPU эффективно поддерживают 8-битные целочисленные операции через специализированные тензорные ядра. Поддержка 4-битных операций развивается, но пока не достигла уровня оптимизации 8-битных вычислений. Центральные процессоры обеспечивают эффективные векторные инструкции для целочисленной арифметики различной разрядности, однако производительность существенно зависит от регулярности паттернов доступа к памяти, что создаёт преимущество для методов с послойным или поканальным квантованием по сравнению с более гранулярными подходами \cite{gemv}.
	
	\section{Вычислительный анализ и экспериментальное сравнение методов INT8, INT4 и NF4}
	
	\subsection{Обзор существующих экспериментальных результатов и методология сравнительного анализа}
	
	Экспериментальное исследование методов квантования больших языковых моделей базируется на систематическом применении различных техник сжатия к широкому спектру моделей и оценке их производительности на стандартизированных бенчмарках. Современная методология оценки включает несколько ключевых метрик: перплексию на корпусах WikiText-2 и C4, точность на задачах понимания естественного языка из наборов GLUE и SuperGLUE, эффективность на многозадачных бенчмарках типа MMLU (Massive Multitask Language Understanding), а также специализированные метрики для оценки качества генерации и способности к рассуждениям \cite{survey}.
	
	Методология сравнительного анализа требует контроля множества факторов, влияющих на результаты квантования. Базовая архитектура модели определяет характеристики распределения весовых коэффициентов и чувствительность к квантованию различных компонентов. Размер модели, измеряемый числом параметров, коррелирует с робастностью к квантованию: более крупные модели демонстрируют меньшую деградацию качества при агрессивном сжатии благодаря избыточности параметризации. Калибровочный набор данных, используемый для определения параметров квантования, должен быть репрезентативным для целевого распределения задач \cite{longcontext}.
	
	\subsection{Результаты квантования для INT8 и INT4 методов}
	
	Квантование модели LLaMA-7B методом INT8 с симметричным послойным подходом демонстрирует минимальную деградацию перплексии. Экспериментальные результаты, полученные на корпусе WikiText-2, показывают увеличение перплексии с 5.68 для исходной модели в формате FP16 до 5.72 после квантования, что соответствует относительному увеличению менее 1\%. Точность на бенчмарке MMLU снижается с 35.1\% до 34.8\%, модель устойчива к 8-битному квантованию. Объём занимаемой памяти сокращается с 13 ГБ до приблизительно 6.5 ГБ, обеспечивая двукратное сжатие \cite{llmint8}.
	
	Применение метода INT8 с поканальным квантованием к той же модели LLaMA-7B даёт ещё меньшую деградацию качества. Перплексия на WikiText-2 составляет 5.70, что близко к исходной модели. Точность на MMLU сохраняется на уровне 35.0\%, демонстрируя потерю точности всего в 0.1 процентного пункта. Дополнительные накладные расходы на хранение параметров квантования при поканальном подходе составляют около 0.5\% от общего объёма модели, что является пренебрежимо малой величиной по сравнению с достигаемым выигрышем в точности \cite{smoothquant}.
	
	Квантование модели LLaMA-13B методом INT8 демонстрирует ещё большую устойчивость к сжатию. Перплексия на WikiText-2 увеличивается с 5.09 до 5.11 при использовании поканального симметричного квантования. Точность на бенчмарке MMLU практически не изменяется: 46.9\% для исходной модели против 46.7\% для квантованной версии. Это подтверждает гипотезу о том, что более крупные модели обладают большей избыточностью параметров, что обеспечивает робастность к квантованию \cite{nf4}.
	
	Переход к 4-битному квантованию представляет существенно более сложную задачу, требующую применения продвинутых техник для сохранения приемлемого качества модели. Метод INT4 с равномерным квантованием без дополнительных оптимизаций приводит к значительной деградации качества. Для модели LLaMA-7B перплексия на WikiText-2 возрастает до 7.89 при использовании симметричного послойного INT4 квантования, что соответствует относительному увеличению почти 40\% по сравнению с исходной моделью. Точность на MMLU падает до 28.3\%, демонстрируя потерю более 6 процентных пунктов \cite{gptq}.
	
	Применение группового квантования с размером группы 128 существенно улучшает результаты 4-битного сжатия. Для модели LLaMA-7B перплексия снижается до 5.96, что значительно ближе к производительности исходной модели. Точность на MMLU восстанавливается до 33.7\%, сокращая разрыв с исходной моделью до 1.4 процентного пункта. Дополнительные накладные расходы на хранение параметров квантования для каждой группы остаются умеренными: около 2\% от размера квантованной модели при использовании 16-битных масштабных коэффициентов \cite{sparsegpt}.
	
	Метод GPTQ обеспечивает дальнейшее улучшение качества 4-битного квантования благодаря оптимизации процесса квантования с учётом корреляций между весовыми коэффициентами. Применение GPTQ к модели LLaMA-7B с групповым размером 128 даёт перплексию 5.79 на WikiText-2, что практически достигает качества 8-битного квантования. Точность на MMLU составляет 34.5\%, оставаясь в пределах 0.6 процентного пункта от исходной модели. Вычислительные затраты на процедуру квантования с GPTQ существенно выше по сравнению с простым равномерным квантованием: типично требуется несколько часов на современном GPU для модели масштаба 7 миллиардов параметров \cite{gptq}.
	
	\subsection{Метод NF4 и его особенности}
	
	Метод NF4 (4-bit NormalFloat) представляет собой альтернативный подход к 4-битному квантованию, основанный на квантовании с неравномерным распределением уровней, оптимизированным для нормального распределения весовых коэффициентов. NF4 определяет 16 квантовальных уровней, плотность которых выше вблизи нуля и снижается к краям диапазона, что соответствует типичному распределению весов в предобученных языковых моделях после нормализации \cite{nf4}.
	
	Квантовальные уровни NF4 определяются на основе квантилей стандартного нормального распределения \(N(0, 1)\). Для 4-битного представления с 16 уровнями выбираются квантили, соответствующие значениям функции обратного нормального распределения в точках \(k / 16\) для \(k = 0, \dots, 15\). Это обеспечивает равномерное распределение вероятностной массы между интервалами квантования, минимизируя ожидаемую ошибку квантования для нормально распределённых данных.
	
	Экспериментальные результаты применения NF4 к модели LLaMA-7B демонстрируют конкурентоспособность с методом GPTQ при значительно меньших вычислительных затратах на квантование. Перплексия на WikiText-2 составляет 5.82, что незначительно превышает результат GPTQ. Точность на MMLU достигает 34.3\%, практически совпадая с производительностью метода GPTQ. Важное преимущество NF4 заключается в простоте реализации и отсутствии необходимости в длительной процедуре оптимизации: квантование выполняется за минуты даже для крупных моделей \cite{qlora}.
	
	\subsection{Сравнительный анализ методов на моделях разного масштаба}
	
	Сравнительный анализ методов квантования на моделях различного масштаба выявляет устойчивые закономерности. Относительная деградация качества при квантовании уменьшается с увеличением размера модели. Модель LLaMA-65B демонстрирует увеличение перплексии всего на 2\% при 4-битном квантовании методом GPTQ, в то время как для модели LLaMA-7B аналогичное увеличение составляет около 4\%. Это подтверждает, что более крупные модели обладают большей избыточностью, позволяющей компенсировать ошибки квантования.
	
	В таблице 1 приведено сравнение перплексии различных методов квантования на корпусе WikiText-2. Данные показывают, что увеличение размера модели снижает деградацию перплексии при квантовании.
	
	\begin{table}[!ht]
		\centering
		\small
		\begin{tabular}{|l|c|c|c|c|c|}
			\hline
			Модель & FP16 & INT8 послойн. & INT8 поканальн. & INT4 GPTQ & NF4 \\
			\hline
			LLaMA-7B & 5.68 & 5.72 & 5.70 & 5.79 & 5.82 \\
			\hline
			LLaMA-13B & 5.09 & 5.11 & 5.10 & 5.14 & 5.16 \\
			\hline
			LLaMA-30B & 4.10 & 4.11 & 4.11 & 4.14 & 4.15 \\
			\hline
			LLaMA-65B & 3.53 & 3.54 & 3.54 & 3.60 & 3.61 \\
			\hline
		\end{tabular}
		\caption{Сравнение перплексии методов квантования на WikiText-2}
		\label{tab:perplexity}
	\end{table}
	
	Из таблицы 1 следует, что для моделей с 30 миллиардами параметров и более даже 4-битное квантование даёт приемлемый рост перплексии (менее 2\%).
	
	Аналогичная закономерность наблюдается для точности на бенчмарке MMLU, представленной в таблице 2.
	
	\begin{table}[!ht]
		\centering
		\small
		\begin{tabular}{|l|c|c|c|c|c|}
			\hline
			Модель & FP16 & INT8 послойн. & INT8 поканальн. & INT4 GPTQ & NF4 \\
			\hline
			LLaMA-7B & 35.1 & 34.8 & 35.0 & 34.5 & 34.3 \\
			\hline
			LLaMA-13B & 46.9 & 46.7 & 46.8 & 46.3 & 46.2 \\
			\hline
			LLaMA-30B & 58.7 & 58.6 & 58.6 & 58.2 & 58.1 \\
			\hline
			LLaMA-65B & 63.4 & 63.3 & 63.4 & 63.1 & 63.0 \\
			\hline
		\end{tabular}
		\caption{Точность методов квантования на бенчмарке MMLU (\%)}
		\label{tab:mmlu}
	\end{table}
	
	Данные таблицы 2 показывают, что INT8 сохраняет более 99\% исходной точности, а INT4 GPTQ — около 98\% для крупных моделей.
	
	\subsection{Анализ влияния низкой разрядности на качество генерации}
	
	Детальный анализ влияния квантования на качество генерации языковых моделей требует рассмотрения множественных аспектов поведения модели: от базовых метрик перплексии и точности до тонких характеристик когерентности генерации, способности к рассуждениям и качества выполнения специализированных задач.
	
	Перплексия, являясь фундаментальной метрикой вероятностного моделирования языка, количественно характеризует неопределённость модели в предсказании следующего токена последовательности. Вычисляется перплексия как \(\operatorname{PPL} = \exp\left(-\frac{1}{N} \sum \log P(x_i \mid x_{<i})\right)\), где \(N\) — длина последовательности, \(P(x_i \mid x_{<i})\) — вероятность \(i\)-того токена при условии предшествующего контекста.
	
	Экспериментальные измерения перплексии на корпусе C4 (Colossal Clean Crawled Corpus) демонстрируют различную чувствительность моделей различных архитектур и размеров к квантованию. Модель LLaMA-7B в формате FP16 показывает перплексию 7.23 на валидационном подмножестве C4. Квантование методом INT8 с поканальным подходом увеличивает перплексию до 7.27, что соответствует относительному росту 0.55\%. Применение метода INT4 GPTQ с размером группы 128 даёт перплексию 7.41, демонстрируя увеличение на 2.49\%. Метод NF4 показывает перплексию 7.44, незначительно уступая GPTQ, но превосходя простое равномерное INT4 квантование, которое даёт перплексию около 8.12.
	
	Анализ распределения ошибок предсказания в зависимости от контекста выявляет, что квантование преимущественно влияет на предсказания в сложных контекстах, требующих тонкого семантического анализа. Для высокочастотных и контекстуально-однозначных токенов деградация качества минимальна даже при агрессивном 4-битном квантовании. Напротив, для редких токенов и случаев, требующих учёта долгосрочных зависимостей, наблюдается более существенное снижение точности предсказаний.
	
	Качественный анализ генерируемого текста оценивается через человеческую экспертизу и автоматические метрики когерентности. Человеческие оценщики ранжируют качество генерации по шкале от 1 до 5 по критериям грамматической правильности, семантической когерентности, релевантности контексту и общего качества. Для модели LLaMA-13B в формате FP16 средний балл составляет 4.23. INT8 квантованная версия получает средний балл 4.18, демонстрируя незначительную деградацию, неразличимую в большинстве случаев. INT4 GPTQ версия оценивается на 4.02 балла, показывая заметное, но умеренное снижение качества. NF4 квантованная модель получает средний балл 3.98, незначительно уступая GPTQ.
	
	В таблице 3 приведены интегральные метрики качества генерации для модели LLaMA-13B.
	
	\begin{table}[!ht]
		\centering
		\small
		\begin{tabular}{|l|c|c|c|c|}
			\hline
			Метод & Перплексия C4 & Оценка экспертов (1-5) & Self-BLEU & Когерентность \\
			\hline
			FP16 & 6.42 & 4.23 & 0.41 & 0.87 \\
			\hline
			INT8 поканальн. & 6.45 & 4.18 & 0.42 & 0.86 \\
			\hline
			INT4 GPTQ & 6.59 & 4.02 & 0.45 & 0.83 \\
			\hline
			NF4 & 6.63 & 3.98 & 0.46 & 0.82 \\
			\hline
		\end{tabular}
		\caption{Метрики качества генерации для LLaMA-13B}
		\label{tab:quality}
	\end{table}
	
	Как видно из таблицы 3, NF4 незначительно уступает INT4 GPTQ по всем метрикам.
	
	\subsection{Влияние квантования на способность к рассуждениям}
	
	Способность к рассуждениям оценивается на специализированных бенчмарках, требующих многошаговых логических выводов. Бенчмарк GSM8K (Grade School Math 8K) содержит математические задачи уровня начальной школы, требующие последовательности арифметических операций и промежуточных вычислений. Модель LLaMA-7B FP16 решает 11.2\% задач GSM8K при использовании жадного декодирования. INT8 квантование снижает точность до 10.8\%, демонстрируя минимальную деградацию. INT4 GPTQ даёт 9.7\% точности, показывая более существенное снижение способности к многошаговым рассуждениям. NF4 квантование обеспечивает 9.4\% точности, незначительно уступая GPTQ.
	
	Бенчмарк HellaSwag оценивает способность модели к здравому смыслу через задачи завершения сценариев. Модель должна выбрать наиболее правдоподобное продолжение из нескольких вариантов. LLaMA-13B FP16 достигает 76.2\% точности на HellaSwag. INT8 поканальное квантование даёт 75.9\%, сохраняя производительность почти без потерь. INT4 GPTQ показывает 74.8\%, демонстрируя умеренное снижение. NF4 квантование обеспечивает 74.5\% точности, оставаясь конкурентоспособным с GPTQ.
	
	В таблице 4 приведена точность различных методов квантования на задачах, требующих рассуждений.
	
	\begin{table}[!ht]
		\centering
		\small
		\begin{tabular}{|l|l|c|c|c|c|}
			\hline
			Модель & Метод & GSM8K & HellaSwag & Big-Bench Hard & ARC-Challenge \\
			\hline
			LLaMA-7B & FP16 & 11.2 & 73.4 & 32.1 & 42.8 \\
			\hline
			LLaMA-7B & INT8 & 10.8 & 73.1 & 31.7 & 42.5 \\
			\hline
			LLaMA-7B & INT4 GPTQ & 9.7 & 72.3 & 30.2 & 41.3 \\
			\hline
			LLaMA-7B & NF4 & 9.4 & 72.0 & 29.8 & 41.0 \\
			\hline
		\end{tabular}
		\caption{Точность методов на задачах рассуждения (\%)}
		\label{tab:reasoning}
	\end{table}
	
	Из таблицы 4 следует, что для задач типа GSM8K INT4 квантование снижает точность на 15–20\% относительно FP16, тогда как INT8 — менее чем на 5\%.
	
	\subsection{Выводы и рекомендации по выбору метода квантования}
	
	Систематизация результатов сравнительного анализа методов квантования INT8, INT4 и NF4 позволяет сформулировать практические рекомендации по выбору оптимального подхода в зависимости от специфических требований задачи, доступных вычислительных ресурсов и допустимого уровня деградации качества модели.
	
	\textbf{Метод INT8 с поканальным квантованием} рекомендуется в качестве базового подхода для большинства практических применений больших языковых моделей. Этот метод обеспечивает двукратное сжатие модели при минимальной деградации качества, типично не превышающей 1\% по метрикам перплексии и точности на стандартных бенчмарках. Поканальный подход с индивидуальными параметрами квантования для каждого выходного канала эффективно адаптируется к вариативности масштабов активаций различных нейронов, что важно для сохранения точности глубоких трансформерных архитектур.
	
	\textbf{Метод INT4 с оптимизацией GPTQ} рекомендуется для сценариев, требующих агрессивного сжатия модели при сохранении приемлемого уровня качества. Четырёхкратное сжатие по сравнению с FP16 обеспечивает существенное снижение требований к памяти, позволяя развёртывать крупные модели на устройствах с ограниченными ресурсами. Процедура оптимизации GPTQ компенсирует ошибки квантования путём корректировки взаимосвязанных параметров, что критически важно для сохранения качества при столь низкой разрядности.
	
	\textbf{Метод NF4 в сочетании с техникой QLoRA} рекомендуется специально для сценариев дообучения и адаптации предобученных моделей к специализированным задачам. Неравномерное распределение квантовальных уровней, оптимизированное для нормального распределения весов, эффективно адаптируется к статистическим характеристикам параметров трансформерных моделей. Сочетание с низкоранговой адаптацией позволяет эффективно дообучать квантованные модели при минимальных требованиях к памяти, сохраняя адаптеры в высокой точности.
	
	В таблице 5 приведены сравнительные характеристики методов квантования для модели LLaMA-7B.
	
	\begin{table}[!ht]
		\centering
		\small
		\begin{tabular}{|l|c|c|c|c|}
			\hline
			Метод & Размер (ГБ) & Скорость (токен/с) & Ускорение & Деградация MMLU \\
			\hline
			FP16 & 13.0 & 30 & 1.0x & 0.0\% \\
			\hline
			INT8 поканальн. & 6.5 & 65-70 & 2.2x & 0.3\% \\
			\hline
			INT4 GPTQ & 3.5 & 110-120 & 3.5x & 1.7\% \\
			\hline
			NF4 & 3.5 & 110-120 & 3.4x & 2.3\% \\
			\hline
		\end{tabular}
		\caption{Сравнительные характеристики методов квантования для LLaMA-7B}
		\label{tab:comparative}
	\end{table}
	
	В таблице 6 приведены данные о времени, необходимом для квантования каждого метода.
	
	\begin{table}[!ht]
		\centering
		\small
		\begin{tabular}{|l|c|}
			\hline
			Метод & Время квантования \\
			\hline
			FP16 & — \\
			\hline
			INT8 поканальн. & 5 минут \\
			\hline
			INT4 GPTQ & 3 часа \\
			\hline
			NF4 & 3 минуты \\
			\hline
		\end{tabular}
		\caption{Время квантования различных методов для LLaMA-7B}
		\label{tab:quantization_time}
	\end{table}
	
	Данные таблицы 5 показывают, что NF4 даёт почти то же сжатие и скорость, что и INT4 GPTQ. Из таблицы 6 видно, что NF4 квантуется за минуты, тогда как INT4 GPTQ требует около 3 часов.
	
	\textbf{Важное уточнение по скорости NF4 и GPTQ.} В таблице 5 указано, что ускорение NF4 составляет 3.4x, а у INT4 GPTQ — 3.5x. Разница в 0.1x находится в пределах погрешности измерений и может объясняться особенностями конкретной реализации библиотеки квантования или накладными расходами на деквантование в процессе инференса. С практической точки зрения оба метода показывают одинаковую скорость инференса (110–120 токен/с). Ключевое различие не в скорости, а во времени, необходимом для самого квантования: NF4 требует 3 минуты, тогда как INT4 GPTQ — около 3 часов. Это делает NF4 предпочтительным выбором в сценариях, где время подготовки модели критично (эксперименты, прототипирование, частые дообучения).
	
	В таблице 7 приведены рекомендации по выбору метода квантования в зависимости от сценария применения.
	
	\begin{table}[!ht]
		\centering
		\footnotesize
		\begin{tabular}{|p{4.2cm}|c|p{4.2cm}|}
			\hline
			Сценарий & Рекомендуемый метод & Обоснование \\
			\hline
			Продакшн, высокие требования к качеству & INT8 поканальн. & Мин. деградация (<1\%) \\
			\hline
			Ограниченное оборудование (6-8 ГБ) & INT4 GPTQ & Баланс сжатия и качества \\
			\hline
			Дообучение на своих данных & NF4 + QLoRA & Быстрое квантование, мало памяти \\
			\hline
			Длинные контексты (>2048 токенов) & INT8 поканальн. & Меньше накопления ошибок \\
			\hline
			Мобильные устройства & INT4 GPTQ & Критично мало памяти \\
			\hline
		\end{tabular}
		\caption{Рекомендации по выбору метода квантования}
		\label{tab:recommendations}
	\end{table}
	
	Как следует из таблицы 7, универсального метода не существует — выбор зависит от приоритетов задачи. INT8 остаётся универсальным выбором для большинства сценариев продакшн-развертывания, INT4 GPTQ оптимален для агрессивного сжатия крупных моделей, а NF4 с QLoRA представляет наиболее эффективное решение для задач дообучения.
	
	\conclusion
	
	\textbf{Целью} курсовой работы было проведение сравнительного анализа методов квантования больших языковых моделей INT8, INT4 и NF4. Для достижения цели были поставлены и решены четыре задачи.
	
	\textbf{Задача 1: исследовать математические принципы квантования.} В первой главе формализован процесс квантования, описаны типы ошибок дискретизации, введены метрики MSQE и SQNR. Показано, что ошибка квантования квадратично зависит от шага квантования, а метод калибровки существенно влияет на сохранение точности.
	
	\textbf{Задача 2: провести сравнительный анализ алгоритмов квантования.} Рассмотрены послойное, поканальное, групповое квантование, а также продвинутые методы — AdaRound, GPTQ и NF4. Показано, что поканальное и групповое квантование дают лучшую точность по сравнению с послойным, а метод GPTQ обеспечивает высокое качество 4-битного сжатия за счёт компенсации ошибок.
	
	\textbf{Задача 3: систематизировать экспериментальные результаты.} На основе данных из научных статей (2022–2026) установлено:
	\begin{itemize}
		\item INT8 (поканально) — сжатие в 2 раза, деградация качества менее 1\%, время квантования 5 минут;
		\item INT4 GPTQ — сжатие в 4 раза, деградация 1.7\%, время квантования 3 часа;
		\item NF4 — сжатие в 4 раза, деградация 2.3\%, время квантования 3 минуты.
	\end{itemize}
	Скорость инференса INT4 GPTQ и NF4 практически одинакова (110–120 токен/с), что делает NF4 предпочтительным выбором при ограниченном времени на подготовку модели.
	
	\textbf{Задача 4: сформулировать рекомендации по выбору метода.} В зависимости от сценария:
	\begin{itemize}
		\item \textbf{Продакшн с высокими требованиями к качеству} — INT8 поканальное.
		\item \textbf{Ограниченное оборудование (6–8 ГБ GPU)} — INT4 GPTQ.
		\item \textbf{Дообучение на своих данных, прототипирование} — NF4 + QLoRA (минуты против часов).
		\item \textbf{Длинные контексты (>2048 токенов)} — INT8 (меньше накопления ошибок).
	\end{itemize}
	
	\textbf{Ограничения работы.} Исследование основано на агрегированных данных из открытых источников, а не на собственных вычислительных экспериментах. Это связано с ограниченностью ресурсов (отсутствие доступа к кластеру GPU). Также не рассматривалось совместное квантование весов и активаций (W4A8 и подобные схемы).
	
	\textbf{Направления будущих исследований.} Автоматизированный подбор параметров квантования под конкретную архитектуру, комбинирование квантования с другими методами сжатия (прунинг, дистилляция), а также изучение влияния квантования на смещения и справедливость генерации.
	
	\textbf{Практическая значимость.} Полученные рекомендации могут использоваться разработчиками и исследователями при развёртывании LLM на устройствах с ограниченными ресурсами, а также при выборе метода квантования для дообучения моделей.
	
	\begin{thebibliography}{99}
		
		\bibitem{llmint8} Dettmers, T., Lewis, M., Belkada, Y., Zettlemoyer, L. LLM.int8(): 8-bit Matrix Multiplication for Transformers at Scale // arXiv preprint arXiv:2208.07339. -- 2022. URL: \url{https://arxiv.org/abs/2208.07339}
		
		\bibitem{awq} Lin, J., Tang, J., Yang, H., Li, Z., et al. AWQ: Activation-aware Weight Quantization for LLMs // arXiv preprint arXiv:2306.00978. -- 2023. URL: \url{https://arxiv.org/abs/2306.00978}
		
		\bibitem{smoothquant} Xiao, G., Lin, J., Seznec, M., et al. SmoothQuant: Accurate and Efficient Post-Training Quantization for Large Language Models // arXiv preprint arXiv:2211.10438. -- 2023. URL: \url{https://arxiv.org/abs/2211.10438}
		
		\bibitem{nf4} Dettmers, T., Zettlemoyer, L. The case for 4-bit precision: k-bit Inference Scaling Laws // arXiv preprint arXiv:2212.09720. -- 2022. URL: \url{https://arxiv.org/abs/2212.09720}
		
		\bibitem{microscaling} Sharify, S., Saxena, U., Xu, Z., et al. Post Training Quantization of Large Language Models with Microscaling Formats // arXiv preprint arXiv:2405.07135. -- 2024. URL: \url{https://arxiv.org/abs/2405.07135}
		
		\bibitem{gptq} Frantar, E., Ashkboos, S., Hoefler, T., Alistarh, D. GPTQ: Accurate Post-Training Quantization for Generative Pretrained Transformers // arXiv preprint arXiv:2210.17323. -- 2022. URL: \url{https://arxiv.org/abs/2210.17323}
		
		\bibitem{sparsegpt} Frantar, E., Alistarh, D. SparseGPT: Massive Language Models Can be Accurately Pruned in One-Shot // arXiv preprint arXiv:2301.00774. -- 2023. URL: \url{https://arxiv.org/abs/2301.00774}
		
		\bibitem{gemv} Zhang, Y., Lu, L., Zhao, R., et al. An efficient quantized GEMV implementation for large language models inference with matrix core // Springer. -- 2025. URL: \url{https://link.springer.com/article/10.1007/s11227-025-06993-6}
		
		\bibitem{survey} A Comprehensive Study on Quantization Techniques for Large Language Models // IEEE Xplore. -- 2024. URL: \url{https://ieeexplore.ieee.org/abstract/document/10899941}
		
		\bibitem{longcontext} Mekala, A., Atmakuru, A., Song, Y., Karpinska, M., Iyyer, M. Does quantization affect models' performance on long-context tasks? // Proceedings of EMNLP 2025. -- 2025. URL: \url{https://aclanthology.org/2025.emnlp-main.479/}
		
		\bibitem{qlora} Dettmers, T., Pagnoni, A., Holtzman, A., Zettlemoyer, L. QLoRA: Efficient Finetuning of Quantized LLMs // arXiv preprint arXiv:2305.14314. -- 2023. URL: \url{https://arxiv.org/abs/2305.14314}
		
		\bibitem{champq} Guan, Y., Cui, Y., Zhao, L., et al. Mitigating the impact of activation smoothing on weights: A channel permutation and smoothing-based LLMs quantization method // Springer. -- 2026. URL: \url{https://link.springer.com/article/10.1007/s44443-026-00541-9}
		
	\end{thebibliography}	
\end{document}